\documentclass[10pt,a4paper]{article}
\usepackage[margin=2.2cm]{geometry}
\usepackage{fontspec}
\setmonofont{DejaVu Sans Mono}[Scale=0.84]
\newfontfamily\unifont{DejaVu Serif}[Scale=0.9]
\usepackage{amsmath,amssymb,mathtools}
\setcounter{MaxMatrixCols}{20}
\usepackage{longtable,booktabs,array}
\usepackage{fancyvrb}
\usepackage[dvipsnames]{xcolor}
\usepackage[colorlinks=true,linkcolor=NavyBlue,urlcolor=NavyBlue]{hyperref}
\usepackage{microtype}
\usepackage{needspace}
\usepackage{tocloft}\setlength{\cftbeforesecskip}{0.45em}
\setlength{\parindent}{0pt}\setlength{\parskip}{0.55em}
\setlength{\emergencystretch}{8em}
\makeatletter
\renewcommand\section{\@startsection{section}{1}{\z@}{-3.5ex \@plus -1ex \@minus -.2ex}{2.3ex \@plus.2ex}{\normalfont\Large\bfseries\raggedright}}
\renewcommand\subsection{\@startsection{subsection}{2}{\z@}{-3.25ex\@plus -1ex \@minus -.2ex}{1.5ex \@plus .2ex}{\normalfont\large\bfseries\raggedright}}
\renewcommand\subsubsection{\@startsection{subsubsection}{3}{\z@}{-3.25ex\@plus -1ex \@minus -.2ex}{1.5ex \@plus .2ex}{\normalfont\normalsize\bfseries\raggedright}}
\makeatother
\renewcommand{\arraystretch}{1.15}
\allowdisplaybreaks
\begin{document}
\title{7+1 repository after the clean-up of 28 October 2026\\[0.4em]\large The state of github.com/once-ere/7plus1 at commit 102f399 — everything that the spreadsheet 7plus1\_status\_2026-10-28.xlsx contains, with every formula it uses}
\author{Repository: github.com/once-ere/7plus1 — prepared for Nash\\ Commit described: 102f399 of 28 October 2026, 13:15 PDT (branch main)\\ Spreadsheet: 7plus1\_status\_2026-10-28.xlsx\\ Data file: tools/status\_2026-10-28.json\\ Document generated by tools/make\_docs.py on 2026-10-28}
\date{}
\maketitle
\begin{abstract}
On 28 October 2026 the documentation of the 7+1 project was completed with the auxiliary torus, the world brane and the radion (sections 25 and 26 of the notebook 7+1.nb), and every loose line of its 30 PDF documents was removed. This document records the state of the repository after that work, at commit 102f399: the 4 commits of that day up to it, the 30 documents with their 478 pages and their loose lines (40, then 12, now 0), the 16 paragraphs that were rewritten, with their formulas, the results of the twelve check programs (504 True, 0 False) and of the 479 checks of the notebook (479 True), the comparison of a fresh clone with the commit, and the private files that are kept out of the repository. It contains the six sheets of the spreadsheet 7plus1\_status\_2026-10-28.xlsx as tables and its 35 formulas with their values; every number comes from tools/status\_2026-10-28.json, the data file of the spreadsheet; the values that LibreOffice computed and stored in the spreadsheet equal those computed here (35 of 35). The conventions of the 7+1 project, TeX’s measure of a loose line and the functions of the spreadsheet are defined here; every equation used is written in this document.
\end{abstract}
\tableofcontents
\Needspace{12\baselineskip}\section*{1. Results at a glance}\addcontentsline{toc}{section}{1. Results at a glance}
The first sheet of the spreadsheet, Summary. Every value is a formula that reads the other sheets (section 8); the last column names the sheet, and the section of this document that reproduces it.

{\footnotesize
\begin{longtable}{>{\raggedright\arraybackslash}p{\dimexpr 0.5000\linewidth-2\tabcolsep\relax}>{\raggedright\arraybackslash}p{\dimexpr 0.1200\linewidth-2\tabcolsep\relax}>{\raggedright\arraybackslash}p{\dimexpr 0.3800\linewidth-2\tabcolsep\relax}}
\toprule
\textbf{item} & \textbf{value} & \textbf{where it comes from} \\
\midrule\endfirsthead
\toprule
\textbf{item} & \textbf{value} & \textbf{where it comes from} \\
\midrule\endhead
\endfoot
\bottomrule\endlastfoot
Commit described (the latest when the spreadsheet was made) & 102f399 & Commits, last row (section 3) \\
Commits of 28 October 2026 up to the commit described & 4 & Commits (section 3) \\
File changes in these commits & 206 & Commits, total (a file changed twice counts twice) (section 3) \\
Documents (Markdown, LaTeX and PDF each) & 30 & Documents (section 4) \\
PDF pages & 478 & Documents, total (section 4) \\
Loose lines in the first build of 28 October & 40 & Documents, total (section 4) \\
Loose lines after commit 0654dcd & 12 & Documents, total (section 4) \\
Loose lines now & 0 & Documents, total (section 4) \\
Overflowing lines now & 0 & Documents, total (section 4) \\
LaTeX errors now & 0 & Documents, total (section 4) \\
Paragraphs rewritten (display, code block or rewording) & 16 & Fixes (section 5) \\
Check programs & 12 & Checks (section 6) \\
Check programs that exit with 0 & 12 & Checks (section 6) \\
Results reported True by the programs & 504 & Checks, total (section 6) \\
Results reported False by the programs & 0 & Checks, total (section 6) \\
Checks in the notebook 7+1.nb & 479 & Checks, the notebook (section 6) \\
Notebook checks that returned True & 479 & Checks, the notebook (section 6) \\
Notebook message (error) cells & 0 & Checks, the notebook (section 6) \\
Markdown and LaTeX files identical after a fresh rebuild & 60 of 60 & Checks, fresh clone (section 6) \\
PDFs with the same text and page count after a fresh rebuild & 30 of 30 & Checks, fresh clone (section 6) \\
Private files on GitHub & 0 & Private files (section 7) \\*
\end{longtable}}
In words: the 4 commits of 28 October 2026 up to 102f399 made 206 file changes (a file changed in two commits counts twice). The 30 documents — each in Markdown, LaTeX and PDF — have 478 pages; the first full build of the day left 40 loose lines in them, 12 remained after commit 0654dcd, and none remain: no line is loose, none overflows, and LaTeX reports no error. For this, 16 paragraphs were rewritten. The 12 check programs all exit with code 0 and report 504 results True and 0 False; the 479 checks of the notebook 7+1.nb all returned True, with no message (error) cell. A fresh clone of 102f399, rebuilt from scratch, reproduces 60 of 60 Markdown and LaTeX files byte for byte and 30 of 30 PDFs with the same text and page count; no private file is on GitHub.

\Needspace{12\baselineskip}\section*{2. Conventions}\addcontentsline{toc}{section}{2. Conventions}
\Needspace{8\baselineskip}\subsection*{2.1 The 7+1 spacetime}\addcontentsline{toc}{subsection}{2.1 The 7+1 spacetime}
The 7+1 project constructs the octonions and the Clifford algebra of an eight-dimensional spacetime with the coordinates $X=(x_0,x_1,\dots,x_7)$: $x_0$ is time, the direction of the octonion unit $e_0=1$; $x_1,x_2,x_3$ are ordinary space; $x_4,\dots,x_7$ are four auxiliary space-like dimensions. Capital indices $A,B$ run over $0,\dots,7$.

\textbf{The signature: $\{7,1,0\}$ and $\mathrm{diag}(+1,-1,\dots,-1)$.} The 7+1 spacetime has one time-like dimension ($x_0$, the direction of the octonion unit $e_0=1$) and seven space-like dimensions ($x_1,x_2,x_3$ ordinary space, $x_4,\dots,x_7$ auxiliary space). Two different objects describe this, and neither may be replaced by the other. (i) The tensor package xAct, used in section 12 of 7+1.nb exactly as in the author’s notebook Einstein-Lovelock-7+1-Nash.nb, declares the metric with \texttt{DefMetric[\{7, 1, 0\}, g71[-a, -b], CD, …]}: the list $\{7,1,0\}$ declares seven space-like dimensions, one time-like dimension and no degenerate dimension — the numbers of entries $+1$ and $-1$ of $\eta_{71}=\mathrm{diag}(-1,+1,\dots,+1)$ (\texttt{η71}), the same spacetime written in the sign convention of 4+4.nb (space-like $=+1$, time-like $=-1$). (ii) The matrix $\eta=-\eta_{71}=\mathrm{diag}(+1,-1,\dots,-1)$ (\texttt{η17}) is simply the squares of the octonion basis: $e_0^2=+1$ for the unit $e_0=1$ (time) and $e_i^2=-1$ for the seven imaginary units $e_1,\dots,e_7$ (space); more generally it is the real part of the product of two basis units, and it gives the quadratic form of the octonions:
\begin{gather*}
\eta_{AB}=\mathrm{Re}(e_Ae_B),\qquad e_0^2=+1,\qquad e_i^2=-1\ \ (i=1,\dots,7),\\ \mathrm{Re}(x\,x)=x_0^2-x_1^2-\cdots-x_7^2=\eta_{AB}\,x^Ax^B\qquad (x=x^Ae_A\in\mathbb O).
\end{gather*}
$\eta$ gives the signs of the Clifford relation $\Gamma_A\Gamma_B+\Gamma_B\Gamma_A=2\,\eta_{AB}\,1_{16}$ of the real 16×16 Dirac matrices: time squares to $+1_{16}$ and each of the seven space directions to $-1_{16}$, exactly like the octonion units (in the notation $Cl(p,q)$, with $p$ generators squaring to $+1$ and $q$ to $-1$, the algebra $Cl(1,7)\cong\mathbb R(16)$). It also gives the signs of the metric components: $g_{00}>0$ along the time-like direction, negative squares along the space-like directions. Every component computed by 7+1.nb follows from the component matrix of the metric; the declared list fixes for xAct the sign of the determinant of the metric, \texttt{SignDetOfMetric[g71]} $=-1$, which is the same in both sign conventions, since $\det \eta=\det \eta_{71}=-1$. 7+1.nb checks \texttt{SignatureOfMetric[g71]} $=\{7,1,0\}$, \texttt{SignDetOfMetric[g71]} $=-1$, $\mathrm{Re}(e_Ae_B)=\eta_{AB}$ and $\mathrm{Re}(x\,x)=\eta_{AB}x^Ax^B$.

\Needspace{8\baselineskip}\subsection*{2.2 Commits, times and counts}\addcontentsline{toc}{subsection}{2.2 Commits, times and counts}
\begin{itemize}
\item A commit is named by the first seven hexadecimal digits of its 40-digit SHA-1 object name in Git (\texttt{102f399} for \texttt{102f399d\allowbreak{}53720ba9\allowbreak{}b6bc05c6\allowbreak{}dbf8d70c\allowbreak{}68c9ac30}); section 3 lists the full names.
\item Times are Pacific Daylight Time (PDT, UTC−7), the time zone of the author’s computer, on which the commits were made; dates are written year-month-day.
\item For a commit, “files changed”, “lines added” and “lines removed” are the three numbers of \texttt{git show --shortstat}: a file counts once per commit that changes it, and a binary file (a PDF) counts as changed but adds and removes no lines.
\item A document is one of the 30 documentation files of the repository at commit 102f399; each exists as Markdown (\texttt{.md}), LaTeX (\texttt{.tex}) and PDF (\texttt{.pdf}). Its “kind” is Main (a document), Symbols (its symbol dictionary) or Provenance (its provenance record); “covers” names the sections of 7+1.nb that it describes; its pages are those that \texttt{pdfinfo} reports.
\item A document was “changed by” a commit if the commit changed its Markdown, LaTeX or PDF file.
\end{itemize}
\Needspace{8\baselineskip}\subsection*{2.3 Loose and overflowing lines}\addcontentsline{toc}{subsection}{2.3 Loose and overflowing lines}
TeX sets each line of a paragraph in a box of the line width $w$. The material of the line has a natural width $w_0$, and its spaces can stretch by at most $s$ (the total stretchability) or shrink by at most $z$ (the total shrinkability) without becoming ugly. When a line must stretch by $t=w-w_0>0$, TeX measures how bad it is by its badness
\begin{equation*}
b=\min\Big(10000,\ 100\Big(\frac ts\Big)^3\Big)\qquad(s>0),\qquad b=10000\quad(s=0,\ t>0),
\end{equation*}
rounded to an integer. A LOOSE line — in the log of LaTeX the message \texttt{Underfull \textbackslash{}\allowbreak{}hbox} — is a line whose badness exceeds 1000, the value of the parameter \texttt{\textbackslash{}\allowbreak{}hbadness} that LaTeX sets: its spaces stretch to more than $10^{1/3}\approx2.154$ times their stretchability,
\begin{equation*}
\text{loose:}\quad b>1000\iff\frac ts>10^{1/3}\approx2.154 .
\end{equation*}
An OVERFLOWING line — the message \texttt{Overfull \textbackslash{}\allowbreak{}hbox} — is wider than the line even with all its spaces shrunk, by more than 0.1 pt, the value of the parameter \texttt{\textbackslash{}\allowbreak{}hfuzz}:
\begin{equation*}
\text{overflowing:}\quad w_0-z-w>0.1\ \text{pt}.
\end{equation*}
The numbers of loose and overflowing lines of a document are the numbers of the messages \texttt{Underfull \textbackslash{}\allowbreak{}hbox} and \texttt{Overfull \textbackslash{}\allowbreak{}hbox} in the log of its last LaTeX run (each message is one box: a line of a paragraph or of a table cell); a LaTeX error is a line of that log that begins with “!”. The documents are compiled with \texttt{latexmk -xelatex}.

\Needspace{8\baselineskip}\subsection*{2.4 The results of the check programs}\addcontentsline{toc}{subsection}{2.4 The results of the check programs}
Each check program prints one line per result. The numbers of results reported True and False by a program $P$ are
\begin{align*}
N_{\rm True}(P)&=\#\{\ell:\ \ell\text{ ends in “: True” or in “(True)”}\}, \\
N_{\rm False}(P)&=\#\{\ell:\ \ell\text{ ends in “: False” or in “(False)”}\},
\end{align*}
where $\ell$ runs over the lines of the output of $P$, $\#$ counts the elements of a set, and spaces between the colon and the word, and after the word, are allowed. The exit code is the status with which the program ends; 0 means success.

\Needspace{12\baselineskip}\section*{3. The commits of 28 October 2026}\addcontentsline{toc}{section}{3. The commits of 28 October 2026}
The sheet Commits. The four commits of 28 October 2026 up to 102f399, on the branch main of github.com/once-ere/7plus1:

{\small
\begin{longtable}{>{\raggedright\arraybackslash}p{\dimexpr 0.1400\linewidth-2\tabcolsep\relax}>{\raggedright\arraybackslash}p{\dimexpr 0.3000\linewidth-2\tabcolsep\relax}>{\raggedright\arraybackslash}p{\dimexpr 0.1800\linewidth-2\tabcolsep\relax}>{\raggedright\arraybackslash}p{\dimexpr 0.1800\linewidth-2\tabcolsep\relax}>{\raggedright\arraybackslash}p{\dimexpr 0.2000\linewidth-2\tabcolsep\relax}}
\toprule
\textbf{commit} & \textbf{date and time (PDT)} & \textbf{files changed} & \textbf{lines added} & \textbf{lines removed} \\
\midrule\endfirsthead
\toprule
\textbf{commit} & \textbf{date and time (PDT)} & \textbf{files changed} & \textbf{lines added} & \textbf{lines removed} \\
\midrule\endhead
\endfoot
\bottomrule\endlastfoot
0654dcd & 28 October 2026, 09:51 & 122 & 24,996 & 4,890 \\
54862f0 & 28 October 2026, 10:15 & 12 & 66 & 19 \\
93c500f & 28 October 2026, 12:40 & 44 & 210 & 93 \\
102f399 & 28 October 2026, 13:15 & 28 & 374 & 355 \\
\textbf{Total} &  & \textbf{206} & \textbf{25,646} & \textbf{5,357} \\*
\end{longtable}}
What each commit contains:

\begin{itemize}
\item \textbf{0654dcd} (\texttt{0654dcdd\allowbreak{}c090f571\allowbreak{}8ecd964c\allowbreak{}73dbda27\allowbreak{}1f81757a}), 28 October 2026, 09:51 — \emph{Sections 25-26: the auxiliary torus, the world brane, and the radion on the 4-torus and the 4-sphere.} Sections 25–26 of 7+1.nb and the \texttt{AUXILIARY\_\allowbreak{}TORUS}, \texttt{WORLD\_\allowbreak{}BRANE} and \texttt{RADION} documents (each with its symbols and provenance files).
\item \textbf{54862f0} (\texttt{54862f02\allowbreak{}8eb7d9a3\allowbreak{}a11b9076\allowbreak{}d79ff19c\allowbreak{}f0219f3f}), 28 October 2026, 10:15 — \emph{Documentation tools: keep the evaluation record in \texttt{tools/\allowbreak{}eval\_\allowbreak{}log\_\allowbreak{}7p1.txt} for fresh-clone rebuilds.} The record of the evaluation of 7+1.nb, kept in the repository so that a fresh clone rebuilds the documents identically.
\item \textbf{93c500f} (\texttt{93c500ff\allowbreak{}2ae19888\allowbreak{}9d3d4abd\allowbreak{}4aafad0e\allowbreak{}a80f42de}), 28 October 2026, 12:40 — \emph{Documentation: no loose lines left in any PDF; long file names and hashes may break.} The clean-up of the loose lines of the documents.
\item \textbf{102f399} (\texttt{102f399d\allowbreak{}53720ba9\allowbreak{}b6bc05c6\allowbreak{}dbf8d70c\allowbreak{}68c9ac30}), 28 October 2026, 13:15 — \emph{Documentation: count every check of \texttt{tools/\allowbreak{}ref\_\allowbreak{}coupled.py} and \texttt{tools/\allowbreak{}ref\_\allowbreak{}torus\_\allowbreak{}brane.py}.} Corrected check counts (50 and 38 instead of 49 and 37), found while this spreadsheet was being prepared.
\end{itemize}
Commit 102f399 corrects two check counts that were found while the spreadsheet was being prepared; the spreadsheet and its three documents were added in the next commit.

\Needspace{12\baselineskip}\section*{4. The 30 documents}\addcontentsline{toc}{section}{4. The 30 documents}
The sheet Documents, in two tables. The first gives what each document is; the second gives its loose lines in three builds — the first full build of the documents on 28 October 2026 (before commit 0654dcd), the rebuild of a fresh clone of 0654dcd, and the rebuild of a fresh clone of 102f399 (“now”) — its overflowing lines and LaTeX errors now, and whether the commits 93c500f and 102f399 changed it.

\Needspace{8\baselineskip}\subsection*{4.1 What each document is}\addcontentsline{toc}{subsection}{4.1 What each document is}
{\footnotesize
\begin{longtable}{>{\raggedright\arraybackslash}p{\dimexpr 0.4000\linewidth-2\tabcolsep\relax}>{\raggedright\arraybackslash}p{\dimexpr 0.1500\linewidth-2\tabcolsep\relax}>{\raggedright\arraybackslash}p{\dimexpr 0.3300\linewidth-2\tabcolsep\relax}>{\raggedright\arraybackslash}p{\dimexpr 0.1200\linewidth-2\tabcolsep\relax}}
\toprule
\textbf{document} & \textbf{kind} & \textbf{covers (sections of 7+1.nb)} & \textbf{pages} \\
\midrule\endfirsthead
\toprule
\textbf{document} & \textbf{kind} & \textbf{covers (sections of 7+1.nb)} & \textbf{pages} \\
\midrule\endhead
\endfoot
\bottomrule\endlastfoot
\texttt{README} & Main & 0–27 (the whole notebook) & 32 \\
\texttt{README\_\allowbreak{}SYMBOLS} & Symbols & 0–27 (the whole notebook) & 13 \\
\texttt{README\_\allowbreak{}PROVENANCE} & Provenance & 0–27 (the whole notebook) & 66 \\
\texttt{DIRAC16COMPLEX} & Main & 7–10 & 22 \\
\texttt{DIRAC16COMPLEX\_\allowbreak{}SYMBOLS} & Symbols & 7–10 & 15 \\
\texttt{DIRAC16COMPLEX\_\allowbreak{}PROVENANCE} & Provenance & 7–10 & 27 \\
\texttt{GRAVITY} & Main & 11–14 & 19 \\
\texttt{GRAVITY\_\allowbreak{}SYMBOLS} & Symbols & 11–14 & 8 \\
\texttt{GRAVITY\_\allowbreak{}PROVENANCE} & Provenance & 11–14 & 15 \\
\texttt{KOHN\_\allowbreak{}SHAM\_\allowbreak{}GRAVITY} & Main & 15–16 & 9 \\
\texttt{KOHN\_\allowbreak{}SHAM\_\allowbreak{}GRAVITY\_\allowbreak{}SYMBOLS} & Symbols & 15–16 & 5 \\
\texttt{KOHN\_\allowbreak{}SHAM\_\allowbreak{}GRAVITY\_\allowbreak{}PROVENANCE} & Provenance & 15–16 & 8 \\
\texttt{STANDARD\_\allowbreak{}MODEL} & Main & 17–19 & 11 \\
\texttt{STANDARD\_\allowbreak{}MODEL\_\allowbreak{}SYMBOLS} & Symbols & 17–19 & 6 \\
\texttt{STANDARD\_\allowbreak{}MODEL\_\allowbreak{}PROVENANCE} & Provenance & 17–19 & 8 \\
\texttt{COUPLED\_\allowbreak{}SOLUTIONS} & Main & 20–23 & 19 \\
\texttt{COUPLED\_\allowbreak{}SOLUTIONS\_\allowbreak{}SYMBOLS} & Symbols & 20–23 & 10 \\
\texttt{COUPLED\_\allowbreak{}SOLUTIONS\_\allowbreak{}PROVENANCE} & Provenance & 20–23 & 16 \\
\texttt{DARK\_\allowbreak{}SECTOR} & Main & 24 & 15 \\
\texttt{DARK\_\allowbreak{}SECTOR\_\allowbreak{}SYMBOLS} & Symbols & 24 & 8 \\
\texttt{DARK\_\allowbreak{}SECTOR\_\allowbreak{}PROVENANCE} & Provenance & 24 & 10 \\
\texttt{AUXILIARY\_\allowbreak{}TORUS} & Main & 25, with 14–24 and 26 & 22 \\
\texttt{AUXILIARY\_\allowbreak{}TORUS\_\allowbreak{}SYMBOLS} & Symbols & 25, with 14–24 and 26 & 13 \\
\texttt{AUXILIARY\_\allowbreak{}TORUS\_\allowbreak{}PROVENANCE} & Provenance & 25, with 14–24 and 26 & 13 \\
\texttt{WORLD\_\allowbreak{}BRANE} & Main & 25–26, with 17–24 & 24 \\
\texttt{WORLD\_\allowbreak{}BRANE\_\allowbreak{}SYMBOLS} & Symbols & 25–26, with 17–24 & 11 \\
\texttt{WORLD\_\allowbreak{}BRANE\_\allowbreak{}PROVENANCE} & Provenance & 25–26, with 17–24 & 13 \\
\texttt{RADION} & Main & 26 & 19 \\
\texttt{RADION\_\allowbreak{}SYMBOLS} & Symbols & 26 & 9 \\
\texttt{RADION\_\allowbreak{}PROVENANCE} & Provenance & 26 & 12 \\
\textbf{Total (30 documents)} &  &  & \textbf{478} \\*
\end{longtable}}
\Needspace{8\baselineskip}\subsection*{4.2 Loose lines, overflowing lines and errors}\addcontentsline{toc}{subsection}{4.2 Loose lines, overflowing lines and errors}
{\footnotesize
\begin{longtable}{>{\raggedright\arraybackslash}p{\dimexpr 0.2700\linewidth-2\tabcolsep\relax}>{\raggedright\arraybackslash}p{\dimexpr 0.1000\linewidth-2\tabcolsep\relax}>{\raggedright\arraybackslash}p{\dimexpr 0.1000\linewidth-2\tabcolsep\relax}>{\raggedright\arraybackslash}p{\dimexpr 0.0800\linewidth-2\tabcolsep\relax}>{\raggedright\arraybackslash}p{\dimexpr 0.1300\linewidth-2\tabcolsep\relax}>{\raggedright\arraybackslash}p{\dimexpr 0.0800\linewidth-2\tabcolsep\relax}>{\raggedright\arraybackslash}p{\dimexpr 0.1200\linewidth-2\tabcolsep\relax}>{\raggedright\arraybackslash}p{\dimexpr 0.1200\linewidth-2\tabcolsep\relax}}
\toprule
\textbf{document} & \textbf{loose, first build} & \textbf{loose after 0654dcd} & \textbf{loose now} & \textbf{overflowing now} & \textbf{errors now} & \textbf{changed by 93c500f} & \textbf{changed by 102f399} \\
\midrule\endfirsthead
\toprule
\textbf{document} & \textbf{loose, first build} & \textbf{loose after 0654dcd} & \textbf{loose now} & \textbf{overflowing now} & \textbf{errors now} & \textbf{changed by 93c500f} & \textbf{changed by 102f399} \\
\midrule\endhead
\endfoot
\bottomrule\endlastfoot
\texttt{README} & 4 & 0 & 0 & 0 & 0 & no & yes \\
\texttt{README\_\allowbreak{}SYMBOLS} & 0 & 0 & 0 & 0 & 0 & no & no \\
\texttt{README\_\allowbreak{}PROVENANCE} & 11 & 1 & 0 & 0 & 0 & yes & yes \\
\texttt{DIRAC16COMPLEX} & 2 & 1 & 0 & 0 & 0 & yes & no \\
\texttt{DIRAC16COMPLEX\_\allowbreak{}SYMBOLS} & 0 & 0 & 0 & 0 & 0 & no & no \\
\texttt{DIRAC16COMPLEX\_\allowbreak{}PROVENANCE} & 6 & 1 & 0 & 0 & 0 & yes & yes \\
\texttt{GRAVITY} & 1 & 1 & 0 & 0 & 0 & yes & no \\
\texttt{GRAVITY\_\allowbreak{}SYMBOLS} & 0 & 0 & 0 & 0 & 0 & no & no \\
\texttt{GRAVITY\_\allowbreak{}PROVENANCE} & 4 & 1 & 0 & 0 & 0 & yes & no \\
\texttt{KOHN\_\allowbreak{}SHAM\_\allowbreak{}GRAVITY} & 0 & 0 & 0 & 0 & 0 & no & no \\
\texttt{KOHN\_\allowbreak{}SHAM\_\allowbreak{}GRAVITY\_\allowbreak{}SYMBOLS} & 0 & 0 & 0 & 0 & 0 & no & no \\
\texttt{KOHN\_\allowbreak{}SHAM\_\allowbreak{}GRAVITY\_\allowbreak{}PROVENANCE} & 1 & 1 & 0 & 0 & 0 & yes & no \\
\texttt{STANDARD\_\allowbreak{}MODEL} & 1 & 1 & 0 & 0 & 0 & yes & no \\
\texttt{STANDARD\_\allowbreak{}MODEL\_\allowbreak{}SYMBOLS} & 1 & 1 & 0 & 0 & 0 & yes & no \\
\texttt{STANDARD\_\allowbreak{}MODEL\_\allowbreak{}PROVENANCE} & 1 & 0 & 0 & 0 & 0 & no & no \\
\texttt{COUPLED\_\allowbreak{}SOLUTIONS} & 2 & 1 & 0 & 0 & 0 & yes & yes \\
\texttt{COUPLED\_\allowbreak{}SOLUTIONS\_\allowbreak{}SYMBOLS} & 0 & 0 & 0 & 0 & 0 & no & no \\
\texttt{COUPLED\_\allowbreak{}SOLUTIONS\_\allowbreak{}PROVENANCE} & 3 & 2 & 0 & 0 & 0 & yes & yes \\
\texttt{DARK\_\allowbreak{}SECTOR} & 1 & 1 & 0 & 0 & 0 & yes & no \\
\texttt{DARK\_\allowbreak{}SECTOR\_\allowbreak{}SYMBOLS} & 0 & 0 & 0 & 0 & 0 & no & no \\
\texttt{DARK\_\allowbreak{}SECTOR\_\allowbreak{}PROVENANCE} & 2 & 0 & 0 & 0 & 0 & no & no \\
\texttt{AUXILIARY\_\allowbreak{}TORUS} & 0 & 0 & 0 & 0 & 0 & no & yes \\
\texttt{AUXILIARY\_\allowbreak{}TORUS\_\allowbreak{}SYMBOLS} & 0 & 0 & 0 & 0 & 0 & no & no \\
\texttt{AUXILIARY\_\allowbreak{}TORUS\_\allowbreak{}PROVENANCE} & 0 & 0 & 0 & 0 & 0 & no & yes \\
\texttt{WORLD\_\allowbreak{}BRANE} & 0 & 0 & 0 & 0 & 0 & no & yes \\
\texttt{WORLD\_\allowbreak{}BRANE\_\allowbreak{}SYMBOLS} & 0 & 0 & 0 & 0 & 0 & no & no \\
\texttt{WORLD\_\allowbreak{}BRANE\_\allowbreak{}PROVENANCE} & 0 & 0 & 0 & 0 & 0 & no & yes \\
\texttt{RADION} & 0 & 0 & 0 & 0 & 0 & no & no \\
\texttt{RADION\_\allowbreak{}SYMBOLS} & 0 & 0 & 0 & 0 & 0 & no & no \\
\texttt{RADION\_\allowbreak{}PROVENANCE} & 0 & 0 & 0 & 0 & 0 & no & no \\
\textbf{Total} & \textbf{40} & \textbf{12} & \textbf{0} & \textbf{0} & \textbf{0} & \textbf{11 yes} & \textbf{9 yes} \\*
\end{longtable}}
The loose lines of the first build were all in the README and in the documents of the earlier stages; the documents of the auxiliary torus, the world brane and the radion had none. Commit 0654dcd let code break inside long file names and hashes, set the headings ragged right and rewrote 5 paragraphs, which removed 28 of the 40 loose lines; commit 93c500f rewrote 11 more paragraphs, which removed the other 12 (section 5).

\Needspace{12\baselineskip}\section*{5. The rewritten paragraphs}\addcontentsline{toc}{section}{5. The rewritten paragraphs}
The sheet Fixes. “How”: Line breaking — a change of the typesetting rules; Display — a long inline formula moved into a displayed equation; Code block — a long piece of code moved into a block of its own; Reworded — the sentence rewritten so that TeX can break its lines evenly; Correction — a correction of content. Every rewritten paragraph keeps its content.

{\footnotesize
\begin{longtable}{>{\raggedright\arraybackslash}p{\dimexpr 0.0500\linewidth-2\tabcolsep\relax}>{\raggedright\arraybackslash}p{\dimexpr 0.2100\linewidth-2\tabcolsep\relax}>{\raggedright\arraybackslash}p{\dimexpr 0.2300\linewidth-2\tabcolsep\relax}>{\raggedright\arraybackslash}p{\dimexpr 0.1100\linewidth-2\tabcolsep\relax}>{\raggedright\arraybackslash}p{\dimexpr 0.3000\linewidth-2\tabcolsep\relax}>{\raggedright\arraybackslash}p{\dimexpr 0.1000\linewidth-2\tabcolsep\relax}}
\toprule
\textbf{\#} & \textbf{document} & \textbf{section of the document} & \textbf{how} & \textbf{change} & \textbf{commit} \\
\midrule\endfirsthead
\toprule
\textbf{\#} & \textbf{document} & \textbf{section of the document} & \textbf{how} & \textbf{change} & \textbf{commit} \\
\midrule\endhead
\endfoot
\bottomrule\endlastfoot
1 & all 30 documents & — & Line breaking & Code may break after \texttt{/\allowbreak{}}, \texttt{\_\allowbreak{}} and \texttt{\textbackslash{}\allowbreak{}}, long hashes every 8 characters; headings are set ragged right & 0654dcd \\
2 & \texttt{README} & 12.4 The nine-channel self-coupling and Kohn–Sham theory in gravitational fields & Display & The definitions of $\Theta^K$, $\sigma_r$ and $\Theta_K$ & 0654dcd \\
3 & \texttt{DIRAC16COMPLEX\_\allowbreak{}PROVENANCE} & 3.4 Grassmann algebra, quantization and the canonical anticommutators & Display & The product and the left derivative of the Grassmann algebra & 0654dcd \\
4 & \texttt{GRAVITY\_\allowbreak{}PROVENANCE} & 2.2 The reference notebook Einstein-Lovelock-7+1-Nash.nb & Code block & The xAct definitions \texttt{CD = LC[g]}, \texttt{RiemannCD}, \texttt{RicciCD}, \texttt{RicciScalarCD}, \texttt{EinsteinCD} & 0654dcd \\
5 & \texttt{STANDARD\_\allowbreak{}MODEL\_\allowbreak{}PROVENANCE} & 4. Derivations & Display & The mixing angle in matter (MSW) & 0654dcd \\
6 & \texttt{COUPLED\_\allowbreak{}SOLUTIONS\_\allowbreak{}PROVENANCE} & 4. Derivations & Display & The mode equation of pair creation & 0654dcd \\
7 & \texttt{COUPLED\_\allowbreak{}SOLUTIONS} & 4. The coupled field equations and the energy–momentum tensors & Display & The minisuperspace action and the identities of its variation & 93c500f \\
8 & \texttt{COUPLED\_\allowbreak{}SOLUTIONS\_\allowbreak{}PROVENANCE} & 4. Derivations & Display & The comoving energy of one massless particle; the minisuperspace action and its identities & 93c500f \\
9 & \texttt{DARK\_\allowbreak{}SECTOR} & 6.1 From the solutions of regimes I and II: overproduction & Display & The Einstein-frame mass of a Kaluza–Klein level and its coupling to the radion, $g_n=|dM_E/d\phi|$ & 93c500f \\
10 & \texttt{DIRAC16COMPLEX} & 6.2 Mode expansion and canonical anticommutators & Display & The five momenta of the exact mode sum & 93c500f \\
11 & \texttt{DIRAC16COMPLEX\_\allowbreak{}PROVENANCE} & 3.3 Real plane-wave spinors & Display & The orthogonal matrix $U$ of the spinors and $U^{\mathsf T}hU$ & 93c500f \\
12 & \texttt{GRAVITY} & 9. How to reproduce & Reworded & The generator command and its files & 93c500f \\
13 & \texttt{GRAVITY\_\allowbreak{}PROVENANCE} & 1. The task & Reworded & The item on the frame matrices and the spin connection & 93c500f \\
14 & \texttt{KOHN\_\allowbreak{}SHAM\_\allowbreak{}GRAVITY\_\allowbreak{}PROVENANCE} & 4. Derivations & Display & The completeness of the $\Theta_K$ and the Fierz identity & 93c500f \\
15 & \texttt{README\_\allowbreak{}PROVENANCE} & 6. The second stage: sections 7–10 of 7+1.nb & Display & The notation of sections 7–10 & 93c500f \\
16 & \texttt{STANDARD\_\allowbreak{}MODEL} & 1. Results at a glance & Reworded & The electroweak scale and the tree-level masses & 93c500f \\
17 & \texttt{STANDARD\_\allowbreak{}MODEL\_\allowbreak{}SYMBOLS} & 1. Conventions & Reworded & The matrix notation & 93c500f \\
18 & 8 documents: \texttt{README}, \texttt{README\_\allowbreak{}PROVENANCE}, \texttt{COUPLED\_\allowbreak{}SOLUTIONS}, \texttt{AUXILIARY\_\allowbreak{}TORUS} and \texttt{WORLD\_\allowbreak{}BRANE}, and the provenance files of the last three & — & Correction & Every check of \texttt{tools/\allowbreak{}ref\_\allowbreak{}coupled.py} (50) and \texttt{tools/\allowbreak{}ref\_\allowbreak{}torus\_\allowbreak{}brane.py} (38) is now counted and listed; before, 49 and 37. A ninth document, \texttt{DIRAC16COMPLEX\_\allowbreak{}PROVENANCE}, changed only in the SHA-256 of \texttt{tools/\allowbreak{}doc\_\allowbreak{}stage3\_\allowbreak{}common.py} that it records & 102f399 \\*
\end{longtable}}
The formulas of the rewritten paragraphs, exactly as the documents write them (a final comma, colon or semicolon that continues the sentence of the document is written here as a full stop). Rows 1, 12, 13 and 18 contain no formula.

\Needspace{8\baselineskip}\subsection*{5.1 Row 2, README: The Clifford products \texorpdfstring{$\Theta^K$}{K}, the sign \texorpdfstring{$\sigma_r$}{σr} and the lowered \texorpdfstring{$\Theta_K$}{K}}\addcontentsline{toc}{subsection}{5.1 Row 2, README: The Clifford products \texorpdfstring{$\Theta^K$}{K}, the sign \texorpdfstring{$\sigma_r$}{σr} and the lowered \texorpdfstring{$\Theta_K$}{K}}
For the 256 increasing multi-indices $K=(A_1<\dots<A_r)$, $r=0,\dots,8$, of the eight real 16×16 Dirac matrices $\Gamma^A$:
\begin{equation*}
\Theta^K=i^{\sigma_r}\Gamma^{A_1}\cdots\Gamma^{A_r},\qquad \sigma_r=\frac{r(r-1)}2\bmod2,\qquad \Theta_K=\eta_{A_1A_1}\cdots\eta_{A_rA_r}\Theta^K .
\end{equation*}
The factor $i^{\sigma_r}$ makes each $\bar\Psi\Theta^K\Psi$ real; $\eta_{AA}$ are the diagonal entries of $\eta=\mathrm{diag}(+1,-1,\dots,-1)$; for $r=0$ the product is $1_{16}$.

\Needspace{8\baselineskip}\subsection*{5.2 Row 3, DIRAC16COMPLEX\_PROVENANCE: The product and the left derivative of the Grassmann algebra}\addcontentsline{toc}{subsection}{5.2 Row 3, DIRAC16COMPLEX\_PROVENANCE: The product and the left derivative of the Grassmann algebra}
For increasing index sets $K$, $K'$ of the anticommuting generators $\theta_k$ ($\theta_k\theta_l=-\theta_l\theta_k$), with $\theta_K$ the product of the generators of $K$ in increasing order:
\begin{align*}
\theta_K\theta_{K'}&=\mathrm{sign}(\pi)\,\theta_{\mathrm{sort}(K\cup K')}\qquad(\text{zero if }K\cap K'\neq\emptyset), \\
\frac{\partial^L(\theta_{k_1}\cdots\theta_{k_r})}{\partial\theta_{k_j}}&=(-1)^{j-1}\,\theta_{k_1}\cdots\widehat{\theta_{k_j}}\cdots\theta_{k_r} .
\end{align*}
$\pi$ is the permutation that sorts the sequence $K$ followed by $K'$, $\mathrm{sort}(K\cup K')$ the sorted union, $\partial^L$ the derivative from the left, and the hat marks the factor that is left out.

\Needspace{8\baselineskip}\subsection*{5.3 Row 4, GRAVITY\_PROVENANCE: The xAct definitions of the connection and the curvature tensors}\addcontentsline{toc}{subsection}{5.3 Row 4, GRAVITY\_PROVENANCE: The xAct definitions of the connection and the curvature tensors}
As the author’s notebook Einstein-Lovelock-7+1-Nash.nb writes them, after a component metric $g$ has been assigned (Wolfram Language with the packages xAct, xCoba and xTras):

\begin{Verbatim}[fontsize=\small,frame=single,framesep=2mm]
CD = LC[g]; RiemannCD = Riemann[CD]; RicciCD = Ricci[CD];
RicciScalarCD = RicciScalar[CD]; EinsteinCD = Einstein[CD];
\end{Verbatim}
After \texttt{CD = LC[g]} the abstract covariant derivative \texttt{CD} is the Levi-Civita connection of the component metric $g$, and \texttt{RiemannCD}, \texttt{RicciCD}, \texttt{RicciScalarCD}, \texttt{EinsteinCD} are its Riemann tensor, Ricci tensor, scalar curvature and Einstein tensor.

\Needspace{8\baselineskip}\subsection*{5.4 Row 5, STANDARD\_MODEL\_PROVENANCE: The mixing angle of two neutrino flavours in matter (MSW)}\addcontentsline{toc}{subsection}{5.4 Row 5, STANDARD\_MODEL\_PROVENANCE: The mixing angle of two neutrino flavours in matter (MSW)}
Diagonalizing the two-flavour Hamiltonian in matter gives the mixing angle in matter $\theta_m$:
\begin{equation*}
\frac{\Delta m^2}{4E}\begin{pmatrix}-\cos2\theta&\sin2\theta\\\sin2\theta&\cos2\theta\end{pmatrix}+\mathrm{diag}(V,0)\quad\Longrightarrow\quad \tan2\theta_m=\frac{\sin2\theta}{\cos2\theta-2EV/\Delta m^2} .
\end{equation*}
$\Delta m^2$ is the difference of the squared masses, $E$ the neutrino energy, $\theta$ the mixing angle in vacuum and $V$ the potential of the electron neutrino in matter.

\Needspace{8\baselineskip}\subsection*{5.5 Row 6, COUPLED\_SOLUTIONS\_PROVENANCE: The mode equation of pair creation}\addcontentsline{toc}{subsection}{5.5 Row 6, COUPLED\_SOLUTIONS\_PROVENANCE: The mode equation of pair creation}
The two-component amplitude $c(\tau)$ of a mode of dirac16complex in the Einstein–Kasner solution of the earliest regime obeys
\begin{equation*}
i\,\frac{dc}{d\tau}=\omega\,(\cos\theta\,\sigma_1+\sin\theta\,\sigma_2)\,c,\qquad \omega=K\sqrt{\tau^{-2p_1}+\tau^{-2q_1}},\qquad \tan\theta=\tau^{p_1-q_1} .
\end{equation*}
$\tau$ is the time since the singularity, $K$ the scaling variable of the mode, $p_1=(1+2\sqrt2)/7$ and $q_1=u_1p_1$ with $u_1=-1+1/\sqrt2$ the Kasner exponents of ordinary and auxiliary space, and $\sigma_1$, $\sigma_2$ Pauli matrices.

\Needspace{8\baselineskip}\subsection*{5.6 Row 7, COUPLED\_SOLUTIONS: The minisuperspace action and the identities of its variation}\addcontentsline{toc}{subsection}{5.6 Row 7, COUPLED\_SOLUTIONS: The minisuperspace action and the identities of its variation}
For the metric $ds^2=n^2dx_0^2-a^2(dx_1^2+dx_2^2+dx_3^2)-b^2(dx_4^2+\dots+dx_7^2)$, with $\mathrm{EL}_q$ the Euler–Lagrange expression for the variable $q$, evaluated at $n=1$:
\begin{align*}
S&=-\frac1{2\kappa}\int n\,a^3b^4\sum_k\alpha_kL_k\,dx_0-\int nW\,dx_0, \\
\mathrm{EL}_n(na^3b^4L_k)&=-2a^3b^4E^{(k)0}{}_0,\qquad \mathrm{EL}_a(na^3b^4L_k)=-6a^2b^4E^{(k)1}{}_1, \\
\mathrm{EL}_b(na^3b^4L_k)&=-8a^3b^3E^{(k)4}{}_4 .
\end{align*}
$\kappa=8\pi G_8$, $\alpha_k$ are the Lovelock couplings ($\alpha_0=2\Lambda$, $\alpha_1=1$), $L_k$ the Lovelock densities, $E^{(k)A}{}_B$ the Lovelock tensors and $W(a,b)$ the comoving energy of the matter.

\Needspace{8\baselineskip}\subsection*{5.7 Row 8, COUPLED\_SOLUTIONS\_PROVENANCE: The comoving energy of one massless particle; the minisuperspace action and its identities}\addcontentsline{toc}{subsection}{5.7 Row 8, COUPLED\_SOLUTIONS\_PROVENANCE: The comoving energy of one massless particle; the minisuperspace action and its identities}
For one massless particle with comoving momenta $k_a$ along ordinary space and $k_b$ along the auxiliary space:
\begin{equation*}
W=\omega=\sqrt{k_a^2/a^2+k_b^2/b^2},\qquad \partial_aW=-\frac{k_a^2}{a^3\omega},\qquad \partial_bW=-\frac{k_b^2}{b^3\omega} .
\end{equation*}
\begin{align*}
S&=-\frac1{2\kappa}\int na^3b^4\sum_k\alpha_kL_k\,dx_0-\int nW\,dx_0, \\
\mathrm{EL}_n(na^3b^4L_k)&=-2a^3b^4E^{(k)0}{}_0,\qquad \mathrm{EL}_a(na^3b^4L_k)=-6a^2b^4E^{(k)1}{}_1, \\
\mathrm{EL}_b(na^3b^4L_k)&=-8a^3b^3E^{(k)4}{}_4 .
\end{align*}
The second display is the minisuperspace action with the identities of its variation, the same equations as those of the previous paragraph, except that the provenance file writes $na^3b^4$ without the thin space of $n\,a^3b^4$; $\omega$ is the energy of the particle, and the derivatives $\partial_aW$, $\partial_bW$ enter the variations with respect to $a$ and $b$.

\Needspace{8\baselineskip}\subsection*{5.8 Row 9, DARK\_SECTOR: The Einstein-frame mass of a Kaluza–Klein level and its coupling to the radion}\addcontentsline{toc}{subsection}{5.8 Row 9, DARK\_SECTOR: The Einstein-frame mass of a Kaluza–Klein level and its coupling to the radion}
A level of Kaluza–Klein mass $\mu_n$ of dirac16complex (mass $m$) has, in the Einstein frame, the mass $M_E(\beta)$, with $M_n=M_E(0)$, and hence the coupling $g_n$ to the canonical radion $\phi=\sqrt{12}M_{\rm Pl}\beta$:
\begin{equation*}
M_E(\beta)=e^{-2\beta}\sqrt{m^2+\mu_n^2e^{-2\beta}},\qquad g_n=\Big|\frac{dM_E}{d\phi}\Big|_{\beta=0}=\frac{2M_n+\mu_n^2/M_n}{\sqrt{12}\,M_{\rm Pl}} .
\end{equation*}
$\beta=\ln(b/b_0)$ is the radion and $M_{\rm Pl}$ the reduced Planck mass.

\Needspace{8\baselineskip}\subsection*{5.9 Row 10, DIRAC16COMPLEX: The five momenta of the exact mode sum}\addcontentsline{toc}{subsection}{5.9 Row 10, DIRAC16COMPLEX: The five momenta of the exact mode sum}
7+1.nb sums the mode algebra exactly, with symbolic $\vec x,\vec y$ and $m=1$, for the symmetric set of five momenta
\begin{equation*}
\vec p\in\big\{\pm(2,2,4,0,0,0,0),\ \pm(0,0,0,2,2,0,0),\ \vec 0\big\}\qquad(E=5,5,3,3,1).
\end{equation*}
$E=\sqrt{m^2+|\vec p|^2}$: $\sqrt{1+24}=5$, $\sqrt{1+8}=3$, $\sqrt1=1$.

\Needspace{8\baselineskip}\subsection*{5.10 Row 11, DIRAC16COMPLEX\_PROVENANCE: The orthogonal matrix \texorpdfstring{$U$}{U} of the spinors}\addcontentsline{toc}{subsection}{5.10 Row 11, DIRAC16COMPLEX\_PROVENANCE: The orthogonal matrix \texorpdfstring{$U$}{U} of the spinors}
With the real spinors $u_s(\vec p)$ of positive and $v_s(-\vec p)$ of negative energy ($s=1,\dots,8$, $u_s^{\mathsf T}u_s=2E$) and the real symmetric plane-wave Dirac Hamiltonian $h(\vec p)$:
\begin{equation*}
U=\frac1{\sqrt{2E}}\big(u_1(\vec p),\dots,u_8(\vec p),v_1(-\vec p),\dots,v_8(-\vec p)\big)\qquad\text{satisfies}\qquad U^{\mathsf T}h(\vec p)\,U=\mathrm{diag}(E1_8,-E1_8).
\end{equation*}
$U$ is real and orthogonal, $U^{\mathsf T}U=1_{16}$; $h(\vec p)=p^i\alpha_i+m\Gamma^0$ with $\alpha_i=\Gamma^0\Gamma^i$ ($i=1,\dots,7$, summed).

\Needspace{8\baselineskip}\subsection*{5.11 Row 14, KOHN\_SHAM\_GRAVITY\_PROVENANCE: The completeness of the \texorpdfstring{$\Theta_K$}{K} and the Fierz identity}\addcontentsline{toc}{subsection}{5.11 Row 14, KOHN\_SHAM\_GRAVITY\_PROVENANCE: The completeness of the \texorpdfstring{$\Theta_K$}{K} and the Fierz identity}
The 256 matrices $\Theta_K$ are complete (the first display), and this gives the Fierz identity (the second display) for commuting spinors $a,b,c,d$:
\begin{equation*}
\sum_K(\Theta_K)_{cd}(\Theta^K)_{ab}=16\,\delta_{ad}\delta_{bc},
\end{equation*}
\begin{equation*}
\sum_{K\in r}(a\Theta_Kb)(c\Theta^Kd)=\sum_JC_{rJ}\,(a\Theta_Jd)(c\Theta^Jb),\qquad C_{rJ}=\frac1{256}\sum_{K\in r}\mathrm{tr}\big(\Theta_K\Theta_J\Theta^K\Theta^J\big).
\end{equation*}
Here $(aMb)=\sum_{i,j}a_iM_{ij}b_j$ for a 16×16 matrix $M$; $K\in r$ runs over the multi-indices of rank $r$ and $J$ over all 256; as subscripts in the first display, $a,b,c,d$ are spinor indices $1,\dots,16$, and $\delta$ is the Kronecker delta.

\Needspace{8\baselineskip}\subsection*{5.12 Row 15, README\_PROVENANCE: The notation of sections 7–10}\addcontentsline{toc}{subsection}{5.12 Row 15, README\_PROVENANCE: The notation of sections 7–10}
With $\Psi$ the field dirac16complex:
\begin{gather*}
\Gamma_A=\Gamma^{(1,7)}_A,\qquad \Gamma^{A}=\eta_{(1,7)}^{AA}\Gamma_A,\qquad \Gamma^0=\Gamma_0,\qquad \chi=\chi^{(1,7)}=i\Gamma^{(1,7)}_8,\\ \bar\Psi=\Psi^\dagger\Gamma^0,\qquad \vec x=(x_1,\dots,x_7),\qquad \vec p=(p^1,\dots,p^7).
\end{gather*}
$\Gamma^{(1,7)}_A$ are the eight real 16×16 Dirac matrices of 7+1.nb, $\eta_{(1,7)}=\mathrm{diag}(+1,-1,\dots,-1)$, and $\Gamma^{(1,7)}_8=\Gamma^{(1,7)}_0\cdots\Gamma^{(1,7)}_7$.

\Needspace{8\baselineskip}\subsection*{5.13 Row 16, STANDARD\_MODEL: The electroweak scale and the masses at tree level}\addcontentsline{toc}{subsection}{5.13 Row 16, STANDARD\_MODEL: The electroweak scale and the masses at tree level}
The reworded sentence states:
\begin{equation*}
v=(\sqrt2G_F)^{-1/2}=246.2196402\ \text{GeV},\qquad m_W=\frac{gv}2,\qquad m_Z=\frac{\sqrt{g^2+g'^2}\,v}2,\qquad m_H^2=2\lambda v^2 .
\end{equation*}
$G_F$ is the Fermi constant, $g$ and $g'$ the gauge couplings of $SU(2)_L$ and $U(1)_Y$, $\lambda$ the quartic coupling of the Higgs field; the document writes the three masses inline, as $m_W=gv/2$ and $m_Z=\sqrt{g^2+g'^2}\,v/2$.

\Needspace{8\baselineskip}\subsection*{5.14 Row 17, STANDARD\_MODEL\_SYMBOLS: The matrix notation}\addcontentsline{toc}{subsection}{5.14 Row 17, STANDARD\_MODEL\_SYMBOLS: The matrix notation}
The reworded convention:
\begin{gather*}
1_n=(\delta_{ij})_{i,j=1}^n,\qquad (M^\dagger)_{ij}=\overline{M_{ji}},\qquad (M^*)_{ij}=\overline{M_{ij}},\\ (A\otimes B)_{(i-1)p+k,\,(j-1)q+l}=A_{ij}B_{kl}\ \ (B\ \text{of size}\ p\times q).
\end{gather*}
$1_n$ is the $n\times n$ identity, $M^\dagger$ the conjugate transpose, $M^*$ the complex conjugate (the bar is complex conjugation of a number) and $A\otimes B$ the Kronecker product, in Mathematica \texttt{KroneckerProduct[A, B]}; the document states these in words, the display writes them out.

\Needspace{12\baselineskip}\section*{6. The checks}\addcontentsline{toc}{section}{6. The checks}
\Needspace{8\baselineskip}\subsection*{6.1 The twelve check programs}\addcontentsline{toc}{subsection}{6.1 The twelve check programs}
The sheet Checks, first table: the outputs of the programs in the rebuild of a fresh clone of 102f399 (section 2.4 defines the counts).

{\footnotesize
\begin{longtable}{>{\raggedright\arraybackslash}p{\dimexpr 0.2500\linewidth-2\tabcolsep\relax}>{\raggedright\arraybackslash}p{\dimexpr 0.3600\linewidth-2\tabcolsep\relax}>{\raggedright\arraybackslash}p{\dimexpr 0.0700\linewidth-2\tabcolsep\relax}>{\raggedright\arraybackslash}p{\dimexpr 0.0700\linewidth-2\tabcolsep\relax}>{\raggedright\arraybackslash}p{\dimexpr 0.0700\linewidth-2\tabcolsep\relax}>{\raggedright\arraybackslash}p{\dimexpr 0.1800\linewidth-2\tabcolsep\relax}}
\toprule
\textbf{program} & \textbf{what it checks} & \textbf{exit code} & \textbf{True} & \textbf{False} & \textbf{verdict} \\
\midrule\endfirsthead
\toprule
\textbf{program} & \textbf{what it checks} & \textbf{exit code} & \textbf{True} & \textbf{False} & \textbf{verdict} \\
\midrule\endhead
\endfoot
\bottomrule\endlastfoot
\texttt{tools/\allowbreak{}verify\_\allowbreak{}7p1\_\allowbreak{}nb.py} & The evaluated 7+1.nb: every check output, the printed numbers against the Python programs, the earlier sections against their backups & 0 & 54 & 0 & All cross-checks passed \\
\texttt{tools/\allowbreak{}ref44.py} & Rebuilds the objects of 4+4.nb & 0 & 10 & 0 & All printed results True \\
\texttt{tools/\allowbreak{}ref71.py} & Rebuilds the octonions and the Clifford matrices of 7+1.nb & 0 & 25 & 0 & All printed results True \\
\texttt{tools/\allowbreak{}extra71.py} & Further identities: boosts, rotations, the classification of the nine signatures, real solutions & 0 & 23 & 0 & All printed results True \\
\texttt{tools/\allowbreak{}ref\_\allowbreak{}dirac16.py} & Sections 7–9: dirac16complex, Kohn–Sham, chiral Kohn–Sham & 0 & 90 & 0 & All checks passed \\
\texttt{tools/\allowbreak{}ref\_\allowbreak{}gravity.py} & Sections 11–14: gravity & 0 & 101 & 0 & All checks passed \\
\texttt{tools/\allowbreak{}ref\_\allowbreak{}ks\_\allowbreak{}gravity.py} & Sections 15–16: Kohn–Sham theory in gravitational fields & 0 & 29 & 0 & All checks passed \\
\texttt{tools/\allowbreak{}ref\_\allowbreak{}sm.py} & Sections 17–19: the Standard Model & 0 & 33 & 0 & All checks passed \\
\texttt{tools/\allowbreak{}ref\_\allowbreak{}coupled.py} & Sections 20–23: the coupled field equations and their solutions & 0 & 50 & 0 & All checks passed \\
\texttt{tools/\allowbreak{}ref\_\allowbreak{}dark.py} & Section 24: the dark sector & 0 & 20 & 0 & All checks passed \\
\texttt{tools/\allowbreak{}ref\_\allowbreak{}torus\_\allowbreak{}brane.py} & Section 25: the auxiliary torus and the world brane & 0 & 38 & 0 & All checks passed \\
\texttt{tools/\allowbreak{}ref\_\allowbreak{}radion.py} & Section 26: the radion & 0 & 31 & 0 & All checks passed \\
\textbf{Total (12 programs)} & \textbf{12 of 12 programs exit with code 0} &  & \textbf{504} & \textbf{0} &  \\*
\end{longtable}}
\Needspace{8\baselineskip}\subsection*{6.2 The notebook 7+1.nb}\addcontentsline{toc}{subsection}{6.2 The notebook 7+1.nb}
The notebook contains 479 checks \texttt{check["name", test]}: 443 print their result where they stand, and 36 are made by the function \texttt{frameChecks} of section 13.1 of the notebook (nine for each of four metrics) and printed by its summary table. The notebook was evaluated in a fresh kernel and finished on 2026-10-28 08:57:47 PDT.

{\footnotesize
\begin{longtable}{>{\raggedright\arraybackslash}p{\dimexpr 0.6200\linewidth-2\tabcolsep\relax}>{\raggedright\arraybackslash}p{\dimexpr 0.3800\linewidth-2\tabcolsep\relax}}
\toprule
\textbf{quantity} & \textbf{value} \\
\midrule\endfirsthead
\toprule
\textbf{quantity} & \textbf{value} \\
\midrule\endhead
\endfoot
\bottomrule\endlastfoot
checks in the notebook & 479 \\
checks that returned True & 479 \\
… printed where they stand & 443 \\
… made by \texttt{frameChecks} (section 13.1), printed by the summary table & 36 \\
message (error) cells & 0 \\
input cells & 189 \\
output cells & 554 \\
evaluation time (seconds) & 283.883 \\
Wolfram version & 15.0.1 for Microsoft Windows (64-bit), fresh kernel \\*
\end{longtable}}
\Needspace{8\baselineskip}\subsection*{6.3 A fresh clone of 102f399, rebuilt from scratch}\addcontentsline{toc}{subsection}{6.3 A fresh clone of 102f399, rebuilt from scratch}
The repository was cloned from GitHub into an empty directory, and every document was rebuilt there with \texttt{tools/\allowbreak{}make\_\allowbreak{}docs.py} (which also runs the twelve programs); the rebuilt files were then compared with the committed ones — Markdown and LaTeX byte for byte, PDF by its text (\texttt{pdftotext}) and its page count (\texttt{pdfinfo}), since a PDF also records the time of its creation.

{\footnotesize
\begin{longtable}{>{\raggedright\arraybackslash}p{\dimexpr 0.6200\linewidth-2\tabcolsep\relax}>{\raggedright\arraybackslash}p{\dimexpr 0.3800\linewidth-2\tabcolsep\relax}}
\toprule
\textbf{quantity} & \textbf{value} \\
\midrule\endfirsthead
\toprule
\textbf{quantity} & \textbf{value} \\
\midrule\endhead
\endfoot
\bottomrule\endlastfoot
files tracked in the repository & 184 \\
Markdown and LaTeX files compared & 60 \\
… identical byte for byte & 60 \\
PDFs compared & 30 \\
… with the same text and page count & 30 \\
documents with layout problems (loose or overflowing lines, errors) & 0 \\
other files changed by the rebuild & 0 \\*
\end{longtable}}
\Needspace{12\baselineskip}\section*{7. The private files}\addcontentsline{toc}{section}{7. The private files}
The sheet Private files: files of the author that are in the working folder of the project but are not published. Their names are not given here.

{\footnotesize
\begin{longtable}{>{\raggedright\arraybackslash}p{\dimexpr 0.2200\linewidth-2\tabcolsep\relax}>{\raggedright\arraybackslash}p{\dimexpr 0.3600\linewidth-2\tabcolsep\relax}>{\raggedright\arraybackslash}p{\dimexpr 0.1200\linewidth-2\tabcolsep\relax}>{\raggedright\arraybackslash}p{\dimexpr 0.3000\linewidth-2\tabcolsep\relax}}
\toprule
\textbf{file} & \textbf{what it is} & \textbf{on GitHub?} & \textbf{how it is kept out} \\
\midrule\endfirsthead
\toprule
\textbf{file} & \textbf{what it is} & \textbf{on GitHub?} & \textbf{how it is kept out} \\
\midrule\endhead
\endfoot
\bottomrule\endlastfoot
The radion note (PDF) & A PDF printout of the author’s e-mail of 28 October 2026 on the radion & No & Listed in \texttt{.git/\allowbreak{}info/\allowbreak{}exclude} on the author’s computer \\
The vielbein notes (PDF) & A PDF printout of the author’s private notes on vielbeins in $n$ dimensions & No & Listed in \texttt{.git/\allowbreak{}info/\allowbreak{}exclude} on the author’s computer \\
A third private PDF & Private working material of the author & No & Listed in \texttt{.gitignore} \\
\texttt{.gitignore} & The ignore rules of the repository & No & Lists itself \\*
\end{longtable}}
\texttt{.gitignore} lists the files that Git does not track in every clone of the repository; \texttt{.git/\allowbreak{}info/\allowbreak{}exclude} lists files that Git does not track in one working folder only, and is itself never published. The check: \texttt{git ls-files} of the fresh clone of 102f399 lists none of these files.

\Needspace{12\baselineskip}\section*{8. The formulas of the spreadsheet}\addcontentsline{toc}{section}{8. The formulas of the spreadsheet}
\Needspace{8\baselineskip}\subsection*{8.1 Cells, ranges and functions}\addcontentsline{toc}{subsection}{8.1 Cells, ranges and functions}
A cell is named by its column letter and its row number (\texttt{B5}: column B, row 5); \texttt{Commits!E6} is the cell E6 of the sheet Commits, and a sheet name with a space is quoted (\texttt{'Private files'!C2}). A range \texttt{A2:A5} is the rectangle of cells from A2 to A5. With $R$ a range, $v(c)$ the value of the cell $c$, $x$ a value and $y$ a value, the functions and operators of the spreadsheet are
\begin{align*}
\mathrm{SUM}(R)&=\sum_{c\in R,\ v(c)\ \text{a number}}v(c),\qquad \mathrm{COUNTA}(R)=\#\{c\in R:\ v(c)\ \text{not empty}\}, \\
\mathrm{COUNTIF}(R,x)&=\#\{c\in R:\ v(c)=x\}\quad(\text{texts compared without regard to case}), \\
x+y&=\text{the sum},\qquad x\,\&\,y=\text{the text of }x\text{ followed by the text of }y .
\end{align*}
\Needspace{8\baselineskip}\subsection*{8.2 The 35 formulas and their values}\addcontentsline{toc}{subsection}{8.2 The 35 formulas and their values}
In the order of the sheets. The values are computed here, from the data file, by an evaluator of these functions; the same values are stored in the spreadsheet.

{\footnotesize
\begin{longtable}{>{\raggedright\arraybackslash}p{\dimexpr 0.1300\linewidth-2\tabcolsep\relax}>{\raggedright\arraybackslash}p{\dimexpr 0.0700\linewidth-2\tabcolsep\relax}>{\raggedright\arraybackslash}p{\dimexpr 0.6000\linewidth-2\tabcolsep\relax}>{\raggedright\arraybackslash}p{\dimexpr 0.2000\linewidth-2\tabcolsep\relax}}
\toprule
\textbf{sheet} & \textbf{cell} & \textbf{formula} & \textbf{value} \\
\midrule\endfirsthead
\toprule
\textbf{sheet} & \textbf{cell} & \textbf{formula} & \textbf{value} \\
\midrule\endhead
\endfoot
\bottomrule\endlastfoot
Summary & B5 & \texttt{=Commits!A5} & 102f399 \\
Summary & B6 & \texttt{=COUNTA(Commits!A2:A5)} & 4 \\
Summary & B7 & \texttt{=Commits!E6} & 206 \\
Summary & B8 & \texttt{=COUNTA(Documents!A2:A31)} & 30 \\
Summary & B9 & \texttt{=Documents!D32} & 478 \\
Summary & B10 & \texttt{=Documents!E32} & 40 \\
Summary & B11 & \texttt{=Documents!F32} & 12 \\
Summary & B12 & \texttt{=Documents!G32} & 0 \\
Summary & B13 & \texttt{=Documents!H32} & 0 \\
Summary & B14 & \texttt{=Documents!I32} & 0 \\
Summary & B15 & \texttt{=COUNTIF(Fixes!D2:D19,"Display")+COUNTIF(Fixes!D2:D19,"Code block")+COUNTIF(Fixes!D2:D19,"Reworded")} & 16 \\
Summary & B16 & \texttt{=COUNTA(Checks!A2:A13)} & 12 \\
Summary & B17 & \texttt{=COUNTIF(Checks!C2:C13,0)} & 12 \\
Summary & B18 & \texttt{=Checks!D14} & 504 \\
Summary & B19 & \texttt{=Checks!E14} & 0 \\
Summary & B20 & \texttt{=Checks!B17} & 479 \\
Summary & B21 & \texttt{=Checks!B18} & 479 \\
Summary & B22 & \texttt{=Checks!B21} & 0 \\
Summary & B23 & \texttt{=Checks!B30\&" of "\&Checks!B29} & 60 of 60 \\
Summary & B24 & \texttt{=Checks!B32\&" of "\&Checks!B31} & 30 of 30 \\
Summary & B25 & \texttt{=COUNTIF('Private files'!C2:C5,"Yes")} & 0 \\
Commits & E6 & \texttt{=SUM(E2:E5)} & 206 \\
Commits & F6 & \texttt{=SUM(F2:F5)} & 25,646 \\
Commits & G6 & \texttt{=SUM(G2:G5)} & 5,357 \\
Documents & D32 & \texttt{=SUM(D2:D31)} & 478 \\
Documents & E32 & \texttt{=SUM(E2:E31)} & 40 \\
Documents & F32 & \texttt{=SUM(F2:F31)} & 12 \\
Documents & G32 & \texttt{=SUM(G2:G31)} & 0 \\
Documents & H32 & \texttt{=SUM(H2:H31)} & 0 \\
Documents & I32 & \texttt{=SUM(I2:I31)} & 0 \\
Documents & J32 & \texttt{=COUNTIF(J2:J31,"Yes")} & 11 \\
Documents & K32 & \texttt{=COUNTIF(K2:K31,"Yes")} & 9 \\
Checks & B14 & \texttt{=COUNTIF(C2:C13,0)\&" of 12 programs exit with code 0"} & 12 of 12 programs exit with code 0 \\
Checks & D14 & \texttt{=SUM(D2:D13)} & 504 \\
Checks & E14 & \texttt{=SUM(E2:E13)} & 0 \\*
\end{longtable}}
\textbf{Check.} The spreadsheet stores, for each formula, the value that LibreOffice computed when it saved the file; 35 of the 35 stored values equal the values of the table.

\Needspace{12\baselineskip}\section*{9. How to reproduce}\addcontentsline{toc}{section}{9. How to reproduce}
\textbf{The spreadsheet.} The command below writes \texttt{7plus1\_\allowbreak{}status\_\allowbreak{}2026-10-28.xlsx} from \texttt{tools/\allowbreak{}status\_\allowbreak{}2026-10-28.json} with openpyxl; if LibreOffice is installed, it then opens and saves the file in headless mode, which computes and stores the value of every formula, and it checks that no formula gives an error:

\begin{Verbatim}[fontsize=\small,frame=single,framesep=2mm]
python3 tools/make_status_xlsx.py
\end{Verbatim}
\textbf{The documents.} This document, its symbol dictionary and its provenance record (Markdown, LaTeX, PDF); the command runs the twelve check programs first:

\begin{Verbatim}[fontsize=\small,frame=single,framesep=2mm]
python3 tools/make_docs.py 7plus1_status_2026-10-28 \
    7plus1_status_2026-10-28_SYMBOLS 7plus1_status_2026-10-28_PROVENANCE
\end{Verbatim}
\textbf{The numbers of the data file} come from the repository — the commits (section 3), the changed documents and the pages:

\begin{Verbatim}[fontsize=\small,frame=single,framesep=2mm]
git show --shortstat --format= <commit>
git show --name-only --format= <commit>
git show 102f399:<document>.pdf > x.pdf && pdfinfo x.pdf
\end{Verbatim}
— and from the logs of LaTeX in \texttt{build/\allowbreak{}docs} (the loose lines), \texttt{build/\allowbreak{}docs/\allowbreak{}runs.json} (the outputs of the programs), the evaluation record \texttt{tools/\allowbreak{}eval\_\allowbreak{}log\_\allowbreak{}7p1.txt} and the evaluated 7+1.nb (the notebook), and a fresh clone (section 6.3).

\end{document}
